% Part: many-valued-logic
% Chapter: infinite-valued-logics
% Section: lukasiewicz

\documentclass[../../../include/open-logic-section]{subfiles}

\begin{document}

\olfileid{mvl}{inf}{luk}

\olsection{લુકાશેવિચનું તર્કશાસ્ત્ર}

\begin{editorial}
  અનંતમૂલ્ય લુકાશેવિચ તર્કશાસ્ત્ર અંગેનો આ વિભાગ ટૂંકું
  ``પ્રારૂપ'' માત્ર છે.
\end{editorial}

\begin{defn}\ollabel{def:lukasiewicz} અનંતમૂલ્ય લુકાશેવિચ
તર્કશાસ્ત્ર~$\LogLuk[\infty]$ નીચેની શ્રેણિક વડે વ્યાખ્યાયિત થાય છે:
\begin{enumerate}
  \item $\lnot$, $\land$, $\lor$, $\lif$ ધરાવતી
  પ્રમાણભૂત વિધાનાત્મક ભાષા~$\Lang L_0$.
  \item સત્યમૂલ્યોનો ગણ $V_\infty$.
  \item એકમાત્ર નિર્દિષ્ટ મૂલ્ય $1$ છે; એટલે $V^+ = \{1\}$.
  \item અસત્ય-અચલનું મૂલ્ય શૂન્ય છે. અન્ય સત્યવિધેયો
  નીચેના વિધેયો વડે અપાય છે:
  \begin{align*}
    \tf{\lnot}[\LogLuk](x) & = 1 - x\\
    \tf{\land}[\LogLuk](x,y) & = \min(x,y)\\
    \tf{\lor}[\LogLuk](x,y) & = \max(x,y)\\
    \tf{\lif}[\LogLuk](x,y) & = \min(1,1-(x-y)) = \begin{cases}
      1 & \text{જો } x \le y\\
      1-(x-y) & \text{અન્યથા.}
    \end{cases}
    \end{align*}
\end{enumerate}
$m$-મૂલ્ય લુકાશેવિચ તર્કશાસ્ત્રની વ્યાખ્યા પણ આવી જ છે;
માત્ર $V = V_m$ છે.
\end{defn}
\sourcecorrection{OLMV-013}{સ્થિર અંગ્રેજી મૂળ
અહીં પ્રમાણભૂત ભાષા~$\Lang L_0$ લે છે, જેમાં
અગાઉથી અસત્ય-અચલ સામેલ છે, પરંતુ તેનું સત્યવિધેય
આપતું નથી. અગાઉની ત્રિમૂલ્ય લુકાશેવિચ શ્રેણિક
સાથે એકરૂપતા રાખવા અહીં તેનું મૂલ્ય શૂન્ય
સ્પષ્ટ કર્યું છે.}

\begin{prop}
  \olref[thr][luk]{def:lukasiewicz} વડે વ્યાખ્યાયિત
  $\LogLuk[3]$ એ \olref{def:lukasiewicz} વડે વ્યાખ્યાયિત
  $\LogLuk[3]$ જેવું જ છે.
\end{prop}

\begin{proof}
  \olref[thr][luk]{def:lukasiewicz}માં આપેલી સંયોજકોની
  સત્યતાસારણીઓની સરખામણી \olref{def:lukasiewicz}નાં
  સમીકરણોથી મળતી નીચેની સત્યતાસારણીઓ સાથે કરીને આ જોઈ શકાય છે:
  \begin{center}
    \begin{tabular}{c|c}
      $\tf{\lnot}$ & \\
      \hline
      $1$ & $0$ \\
      $1/2$ & $1/2$ \\
      $0$ & $1$
    \end{tabular}
    \quad
    \begin{tabular}{c|ccc}
      $\tf{\land}[\LogLuk[3]]$ & $1$ & $1/2$ & $0$ \\
      \hline
      $1$ & $1$ & $1/2$ & $0$ \\
      $1/2$ & $1/2$ & $1/2$ & $0$\\
      $0$ & $0$ & $0$ & $0$
    \end{tabular}
    \\[2ex]
    \begin{tabular}{c|ccc}
      $\tf{\lor}[\LogLuk[3]]$ & $1$ & $1/2$ & $0$ \\
      \hline
      $1$ & $1$ & $1$ & $1$ \\
      $1/2$ & $1$ & $1/2$ & $1/2$ \\
      $0$ & $1$ & $1/2$ & $0$
    \end{tabular}
    \quad
    \begin{tabular}{c|ccc}
      $\tf{\lif}[\LogLuk[3]]$ & $1$ & $1/2$ & $0$ \\
      \hline
      $1$ & $1$ & $1/2$ & $0$ \\
      $1/2$ & $1$ & $1$ & $1/2$ \\
      $0$ & $1$ & $1$ & $1$
    \end{tabular}
  \end{center}
\end{proof}

\begin{prop}\ollabel{prop:luk-infty-m}
  જો $\Gamma \Entails[\LogLuk[\infty]] !B$, તો બધાં~$m \ge 2$
  માટે $\Gamma \Entails[\LogLuk[m]] !B$.
\end{prop}

\begin{proof}
  અભ્યાસપ્રશ્ન.
\end{proof}

\begin{prob}
  \olref[mvl][inf][luk]{prop:luk-infty-m} સાબિત કરો.
\end{prob}

હકીકતમાં, સાન્ત આધારગણ માટે વિપરીત દિશા પણ સાચી છે.
\sourcecorrection{OLMV-014}{સ્થિર અંગ્રેજી મૂળ
વિપરીત દિશા કોઈ મર્યાદા વિના કહે છે. સાન્ત
આધારગણ માટે વિરોધી મૂલ્યાંકનમાં આવતા માત્ર
સાન્ત ઘણા સંમેય મૂલ્યોને એક જ સાન્ત જાળીમાં
સમાવી શકાય છે. પરંતુ અનંત આધારગણ જુદા જુદા
છેદવાળાં અસંખ્ય ચલોનાં મૂલ્યો નિર્ધારિત કરી શકે;
બધી સાન્ત જાળીઓમાં તેના આધાર અસંતોષ્ય છતાં
અનંત સંમેય જાળીમાં સંતોષ્ય હોઈ શકે. તેથી અહીં
દાવો સાન્ત આધારગણ સુધી મર્યાદિત કર્યો છે.}

અનંતમૂલ્ય લુકાશેવિચ તર્કશાસ્ત્ર સૌથી પ્રચલિત અસ્પષ્ટ
તર્કશાસ્ત્ર છે. અસ્પષ્ટ તર્કશાસ્ત્રના સાહિત્યમાં પ્રેરણ સંયોજકને
ઘણી વાર $\lnot !A \lor !B$ તરીકે વ્યાખ્યાયિત કરે છે. ત્યારે
અનંતમૂલ્ય પ્રબળ ક્લીની તર્કશાસ્ત્ર મળશે.

\begin{prob}
  $(p \lif q) \lor (q \lif p)$ એ~$\LogLuk[\infty]$ની
  પુનરુક્તિ છે તે બતાવો.
\end{prob}

\end{document}
